\section{Resumen Ejecutivo}

El presente proyecto de investigación aborda una problemática crítica en la agroindustria cafetera colombiana: la limitada capacidad para estandarizar y escalar el proceso de secado de cafés diferenciados (como los procesos honey y natural), lo que genera alta variabilidad en la calidad del producto final y pérdidas económicas para los caficultores.

Actualmente, el secado del café en Colombia se realiza mayoritariamente bajo esquemas tradicionales basados en el control de temperatura y en condiciones ambientales variables, lo cual dificulta el manejo preciso de la humedad del grano, especialmente en procesos no convencionales. Esta situación se ve agravada por las limitaciones técnicas y de infraestructura a nivel de finca, donde muchos productores no cuentan con las herramientas necesarias para garantizar condiciones óptimas de procesamiento.

En este contexto, el proyecto propone el desarrollo y escalamiento de una tecnología de secado de café basada en bomba de calor, cuyo enfoque innovador radica en el control del proceso a partir de la humedad relativa en lugar de la temperatura. Esta estrategia permite un manejo más preciso de la cinética de secado, favoreciendo la conservación de la calidad físico-química y sensorial del grano, así como la reducción de defectos asociados a procesos inadecuados.

La iniciativa se fundamenta en desarrollos previos realizados durante una tesis de pregrado en la Universidad del Valle, donde se evaluó un prototipo piloto. A partir de estos resultados, el presente proyecto busca llevar esta tecnología a un nivel superior mediante su diseño, optimización y escalamiento a capacidades industriales (entre 5 y 20 toneladas por lote), incorporando además elementos complementarios como sistemas de movimiento controlado del producto y la integración de fuentes de energía renovable, específicamente energía solar.

Adicionalmente, el proyecto se articula con transformaciones recientes en el modelo productivo cafetero, como el esquema de compra de café en cereza (CERPER), promovido por Cenicafé y la Federación Nacional de Cafeteros, el cual busca centralizar el procesamiento en plantas especializadas. En este sentido, se plantea como escenario de implementación inicial una planta de beneficio proyectada en el municipio de Sevilla, Valle del Cauca, donde esta tecnología podría contribuir significativamente a mejorar la eficiencia del proceso, garantizar la calidad del producto y generar mayores ingresos para los caficultores.

La ejecución del proyecto se justifica por su potencial para impactar positivamente tanto a nivel técnico como económico y social, al facilitar la industrialización del secado de cafés diferenciados, reducir pérdidas en la cadena productiva y fortalecer la competitividad del café colombiano en mercados internacionales.

Para su desarrollo, se emplearán herramientas de diseño e ingeniería de procesos, modelado y simulación del secado, validación experimental, análisis de calidad del producto, así como evaluaciones de viabilidad económica y de impacto social. De esta manera, se espera generar una solución tecnológica integral que contribuya a la modernización y sostenibilidad del sector cafetero colombiano.